\section{Proposed epistemic reject-option predictor}
\label{sec:ProposedEpistemRejOption}
The goal of our framework is to construct a reject-option predictor whose accept–reject decisions are based solely on epistemic uncertainty. We formulate this predictor as the solution to a well-posed decision-making problem by departing from standard Bayesian learning, which minimizes expected loss~\eqref{equ:BaysianRejOptObj}, and instead minimizing expected regret. Regret is defined as the performance gap between the predictor inferred from data and the Bayes-optimal predictor with full knowledge of the underlying parameters. The resulting predictor abstains whenever the expected regret for a given input exceeds a specified rejection cost, which determines the maximum performance degradation, relative to the optimal predictor, that we are willing to accept.

%The goal of our framework is to construct a reject-option predictor whose accept–reject decisions are based solely on epistemic uncertainty. We formulate this predictor as the solution to a well-posed decision-making problem. We depart from standard Bayesian learning, which minimizes expected loss~\eqref{equ:BaysianRejOptObj}, and instead we minimize expected regret, defined as the performance gap between the predictor inferred from data and the Bayes-optimal predictor with full knowledge of the underlying parameters. The resulting predictor abstains whenever the expected regret for a given input exceeds a specified rejection cost. This rejection cost therefore determines the maximum performance degradation, relative to the optimal predictor, that we are willing to accept.

%The goal of our framework is to construct a reject-option predictor that bases its accept–reject decisions solely on epistemic uncertainty. We define this predictor as the solution to a well-posed decision-making problem. The central idea is to depart from standard BL, which minimizes expected loss~\equ{equ:BaysianRejOptObj}, and instead minimize expected regret, defined as the performance gap between the predictor inferred from data and the Bayes-optimal predictor with full knowledge of the underlying parameters. The predictor abstains when the expected regret for a given input data exceeds a specified rejection cost. The reject cost thus specifies how much performance decrease, relative to the optimal predictor, we are willing to still accept.

Following the BL paradigm, our aim is to construct a reject-option predictor $Q \colon \mathcal{X} \times (\mathcal{X} \times \mathcal{Y})^m \rightarrow \mathcal{Y} \cup \{ \text{reject} \}$,  which, given training data $D \in \mathcal{D}$ and an input $x \in \mathcal{X}$, either predicts $Q(x, D) \in \mathcal{Y}$ or opts to reject, $Q(x, D) = \text{reject}$. Again, we use~\equ{equ:BLRejectOptionPred} to represent $Q$ by a pair $(H,C)$, where  $H\colon \mathcal{X} \times (\mathcal{X} \times \mathcal{Y})^m \rightarrow \mathcal{Y}$ is a predictor and $C\colon \mathcal{X} \times (\mathcal{X} \times \mathcal{Y})^m \rightarrow \{0,1\}$ is a selector.

Let $\ell\colon\SY\times\SY\rightarrow\Re_+$ be a prediction loss and $\delta\in\Re_+$ a rejection cost. Then, the decision of the epistemic reject-option predictor $Q=(H,C)$ on a data sample $(x,D)\sim p(x,D)$ is penalized as follows:

\noindent 1. If it predicts, $C(x,D)=1$, the penalty equals to the expected conditional regret
  \[
    E(x,D) = \EE_{(\theta,y)\sim p(\theta,y|x,D)}\big [ \ell(y,H(x,D)) - \ell(y,h(x,\theta)) \big ]\:,
  \]
  i.e. the expected performance gap between the predictor $H(x,D)$ and the Bayes-optimal predictor $h(x,\theta)$ knowing the true parameter $\theta$.

\noindent 2. If it rejects, $C(x,D)=0$, the penalty equals $\delta$.

The overall performance of the epistemic reject-option predictor $Q=(H,C)$ is then evaluated by:
\begin{multline}
\label{equ:EpistemicRejOptObj}
   R_\delta(H,C)=\\ \EE_{(x,D)\sim p(x,D)} \Big [C(x,D)\, E(x,D)+ \neg C(x,D)\,\delta 
       \Big ]    \:.
\end{multline}

The following theorem characterizes the minimizer of $R_\delta(H,C)$, which we refer to as the optimal epistemic reject-option predictor.
%
\begin{theorem}\label{thm:EpistemicRejOpt} Let $Q_E=(H_*,C_E)$ be a reject-option predictor represented by the optimal Bayesian predictor $H_*$, defined in~\equ{equ:BLOptBayes}, and the selector $C_E(x,D) = \one{ E_*(x,D) \leq \delta }$
where 
\begin{multline}
  \label{equ:ConditionalRegret}
  E_*(x,D) = \\\EE_{(\theta,y)\sim p(\theta,y|x,D)}\big [ \ell(y,H_*(x,D)) - \ell(y,h(x,\theta)) \big ]\:.
\end{multline}
%
Then, $Q_E$ is a minimizer of $R_\delta(H,C)$ defined in~\equ{equ:EpistemicRejOptObj}. 
We refer to \(Q_E\) as the \emph{optimal epistemic reject-option predictor}.

\end{theorem}
%
We defer the proof to Appendix~\ref{appendix:ProofOfEpistemicTheorem}.

The expected conditional regret $E_*(x, D)$ forms a basis of the accept–reject decisions of the epistemic reject-option predictor $Q_E(x, D)$. The value of $E_*(x, D)$ admits a clear interpretation: for a given input $x$ and training set $D$, it quantifies the expected performance gap between the learned predictor and the Bayes-optimal predictor with full knowledge of the data-generating process. This makes $E_*(x, D)$ a natural measure of epistemic uncertainty, and we refer to it as such in the following. Consequently, the epistemic reject-option predictor accepts inputs for which the predictions of the learned predictor are sufficiently close, where closeness is controlled by $\delta$, to those of the theoretically optimal Bayes predictor.